\documentclass{article}
% Estilo NeurIPS oficial: neurips_2026.sty (o 2025) en esta carpeta; si falta, se usa una réplica (neurips_emul.sty).
\IfFileExists{neurips_2026.sty}{\usepackage[preprint]{neurips_2026}}{%
  \IfFileExists{neurips_2025.sty}{\usepackage[preprint]{neurips_2025}}{\usepackage{neurips_emul}}}
% El estilo oficial escribe 'Abstract' fijo: en un informe en español lo titulamos 'Resumen'
\makeatletter\@ifpackageloaded{neurips_emul}{}{\renewenvironment{abstract}%
  {\vskip 0.075in\centerline{\large\bf Resumen}\vspace{0.5ex}\begin{quote}}{\par\end{quote}\vskip 1ex}}\makeatother
\usepackage[utf8]{inputenc}
\usepackage[T1]{fontenc}
% Nombres en español sin depender de babel-spanish (compila igual en Overleaf o TeX Live mínimo)
\AtBeginDocument{\renewcommand{\abstractname}{Resumen}\renewcommand{\figurename}{Figura}%
  \renewcommand{\tablename}{Tabla}\renewcommand{\refname}{Referencias}}
\usepackage{booktabs,graphicx,amsmath,amssymb,xcolor,url,hyperref,microtype,natbib}
\graphicspath{{figures/}}

% --- utilidades: tablas/figuras autogeneradas por el pipeline ---
\newcommand{\pendiente}[1]{\textcolor{red}{[\textbf{Pendiente:} #1]}}
\newcommand{\tabla}[1]{\IfFileExists{tables/#1.tex}{\input{tables/#1.tex}}{\pendiente{tabla \texttt{#1} — se genera con \texttt{run\_all.sh}}}}
\newcommand{\figura}[3]{\begin{figure}[!htbp]\centering
  \IfFileExists{figures/#1.pdf}{\includegraphics[width=#2\linewidth]{#1.pdf}}{\fbox{\pendiente{figura #1}}}
  \caption{#3}\label{fig:#1}\end{figure}}

\title{Detección de fauna silvestre en cámaras trampa con YOLOv8:\\
desbalance, generalización a cámaras nuevas y monitoreo ecológico}

\author{Grupo 4 \\ Martín Alexis Valdiviezo Palma \quad Gabriel Ángel Nizama Duque \\
Juan José Pino Reyes \quad César Luis Jiménez Miranda \\
Maestría en Inteligencia Artificial, FIIS -- Universidad Nacional de Ingeniería}

\begin{document}
\maketitle

\begin{abstract}
Las cámaras trampa generan más imágenes de las que un equipo humano puede revisar, y la mayoría
están vacías. Entrenamos un detector YOLOv8 sobre el benchmark Caltech Camera Traps (CCT-20) para
localizar y clasificar 15 clases. Con la partición
oficial que separa ubicaciones de cámara vistas (\emph{cis}) de ubicaciones nunca vistas (\emph{trans}),
el mAP@0.5 cae de 0.720 a 0.455 (26.5 puntos), el desempeño por cámara nueva varía entre 0.16 y 0.74,
y las especies raras pasan de 0.50 a 0.15. Al borrar al animal por \emph{inpainting}, la misma especie
se sigue detectando en el 16\,\% de las imágenes de cámaras vistas frente al 9\,\% de las nuevas, indicio
de que parte del desempeño en cámaras conocidas proviene del contexto de la escena. Los animales pequeños
son el caso más difícil (recall 0.45 en cámaras nuevas). Un módulo de monitoreo basado en reglas,
aplicado a una cámara no vista con 6\,276 fotos, estima el total de eventos independientes con un error
del 6\,\%, pero ese acierto agregado oculta errores grandes por especie: inventa una clase ausente, omite
dos especies reales y sobreestima el porcentaje de fotos vacías.
\end{abstract}

% =====================================================================
\section{Motivación}
Las cámaras trampa, activadas por movimiento, son hoy la herramienta estándar para inventariar
mamíferos medianos y grandes en áreas protegidas, incluidos programas en la Amazonía y los Andes
peruanos. Su cuello de botella es la revisión: un solo proyecto produce cientos de miles de fotos y,
en CCT, alrededor del 70\,\% son disparos falsos por viento, vegetación o cambios de luz.
Automatizar la revisión tiene dos usos distintos: (i) \emph{filtrar} fotos vacías para ahorrar
tiempo de especialista, y (ii) \emph{identificar y contar} especies para alimentar indicadores
ecológicos. El reto técnico no es solo detectar objetos: las especies siguen una distribución de
cola larga, los animales aparecen pequeños, lejanos, de noche o tapados por vegetación, y, sobre
todo, cada cámara está fija y ve siempre el mismo fondo, lo que invita al modelo a aprender
atajos~\citep{geirhos2020shortcut} que colapsan al instalarlo en un sitio nuevo~\citep{beery2018recognition}.

\paragraph{Contribuciones.} (1) Un pipeline reproducible de un solo comando que convierte CCT-20 a
YOLO con manejo explícito de fotos vacías y de animales sin caja. (2) Un análisis por especie del
efecto del desbalance (503:1) y de su interacción con el cambio de cámara. (3) Una medición de la brecha cis/trans complementada
con dos pruebas de dependencia del fondo (energía Grad-CAM dentro de la caja y enmascarado del
animal). (4) Un análisis por tamaño y oclusión. (5) Un módulo de monitoreo ecológico cuyos
indicadores se validan contra la verdad de campo.

% =====================================================================
\section{Trabajos relacionados}
\textbf{Clasificación de cámaras trampa.} \citet{norouzzadeh2018} mostraron, sobre Snapshot
Serengeti, que una CNN puede identificar, contar y describir animales con exactitud cercana a la de
voluntarios, y \citet{tabak2019} extendieron el enfoque a especies norteamericanas. Ambos trabajos
clasifican la imagen completa, sin localizar al animal.
\textbf{Detección.} \citet{schneider2018} compararon Faster R-CNN y YOLO en datos de cámaras trampa y
encontraron que la detección mejora el conteo y la localización. MegaDetector~\citep{beery2019megadetector}
popularizó un detector genérico \emph{animal/persona/vehículo} seguido de un clasificador de especie;
nuestro enfoque es de una sola etapa con clases de especie, siguiendo la familia YOLO~\citep{redmon2016yolo,jocher2023yolov8,wang2024yolov9}.
\textbf{Generalización geográfica.} \citet{beery2018recognition} introdujeron CCT-20 y sus particiones
\emph{cis/trans} precisamente para cuantificar la caída de desempeño en cámaras no vistas, y
WILDS~\citep{koh2021wilds} incluye iWildCam como caso canónico de desplazamiento de distribución.
\textbf{Desbalance.} Para detección de cola larga, \citet{gupta2019lvis} proponen \emph{Repeat Factor
Sampling} (RFS); la pérdida focal~\citep{lin2017focal} aborda el desbalance fondo/objeto dentro de
cada imagen. \textbf{Interpretabilidad.} Usamos Grad-CAM~\citep{selvaraju2017gradcam}
sobre el backbone del detector. \textbf{Ecología.} Los índices de abundancia relativa con criterio de
independencia de 30 minutos siguen la práctica de \citet{obrien2003}.

% =====================================================================
\section{Datos}
Caltech Camera Traps completo contiene 243\,100 imágenes de 140 cámaras (105\,GB). Trabajamos con su
subconjunto de referencia CCT-20 (LILA BC, $\sim$6\,GB), publicado por los mismos autores para evaluar
generalización entre ubicaciones, con imágenes reducidas a 1024\,px de lado máximo y cajas delimitadoras.
La elección responde a dos razones: (i) el conjunto completo no cabe en el disco ni en la cuota de
almacenamiento del entorno de cómputo disponible (Google Colab con GPU T4), y (ii) la partición que
exige el experimento de generalización está definida sobre las mismas 20 ubicaciones de CCT-20: la partición oficial \texttt{eccv\_train/eccv\_val} de
LILA corresponde a \texttt{train}+\texttt{cis} (ubicaciones vistas) y \texttt{trans} (ubicaciones
nuevas) del artículo original~\citep{beery2018recognition}. Entrenamos con \texttt{train},
seleccionamos modelo y umbral con \texttt{cis\_val}, y reportamos sobre \texttt{cis\_test} y
\texttt{trans\_test}; \texttt{trans\_val} no se usa para ninguna decisión.
La Tabla~\ref{tab:splits} resume la composición y la Tabla~\ref{tab:species} la prevalencia por
especie. Dos rasgos de CCT-20 condicionan el resto del trabajo.
Primero, a diferencia del conjunto completo ($\sim$70\,\% de fotos vacías), tiene pocas fotos vacías:
ninguna en \texttt{train} y alrededor del 7.5\,\% en las particiones de prueba. Segundo, entre el 11\,\%
y el 15\,\% de las imágenes de cada partición tienen etiqueta de especie pero ninguna caja.

\tabla{eda_splits}
\tabla{eda_species}

% =====================================================================
\section{Metodología}
\subsection{Preparación y manejo de fotos vacías}
Las anotaciones en formato COCO Camera Traps se convierten a YOLO (centro, ancho, alto normalizados),
reescalando las cajas al tamaño real del archivo. Tres casos por imagen: (a) \emph{con caja}: se
entrena normalmente; (b) \emph{vacía}: el pipeline la incluye como \textbf{imagen de fondo} con archivo
de etiquetas vacío (fracción configurable, 10\,\% de las imágenes con animal por defecto), de modo que
toda predicción sobre ella es un falso positivo; (c) \emph{animal sin caja}: se excluye del
entrenamiento de detección porque el animal no anotado se aprendería como fondo, pero se conserva para
el módulo de monitoreo. Como \texttt{train} de CCT-20 no contiene fotos vacías, el caso (b) no se activó
y el modelo no vio ninguna foto vacía al entrenar. El fondo se aprende igualmente de forma implícita: en
YOLOv8, todo punto de la grilla no asignado a una caja recibe objetivo cero en la rama de clasificación,
así que el fondo de cada foto con animal actúa como ejemplo negativo. Las fotos vacías de
\texttt{cis\_val} se usan para fijar el umbral operativo y las de prueba para medir falsos positivos.

\subsection{Modelo y entrenamiento}
YOLOv8s~\citep{jocher2023yolov8} preentrenado en COCO~\citep{lin2014coco}, entrada 640\,px, AdamW
($\text{lr}_0=10^{-3}$, coseno), 30 épocas con parada temprana (paciencia 10), lote 16, semilla 42.
Aumento de datos: mosaico (desactivado en las 5 últimas épocas), traslación 0.1, escala 0.5, volteo
horizontal 0.5, jitter HSV con $v=0.4$ para cubrir capturas infrarrojas nocturnas; \emph{sin}
rotación ni volteo vertical, porque las cámaras están fijas y los animales no aparecen invertidos.

\subsection{Mitigación del desbalance de especies}
Con RFS~\citep{gupta2019lvis}, cada especie $c$ con frecuencia de imagen $f_c$ recibe
$r_c=\max(1,\sqrt{t/f_c})$, $t=0.1$; cada imagen se repite $r_i=\max_{c\in i} r_c$ veces
(redondeo estocástico). Así las especies raras se ven más a menudo sin alterar la pérdida.
RFS está implementado en el pipeline (experimento \texttt{rfs}), pero por el límite de cuota de GPU
no alcanzamos a entrenarlo: todos los resultados corresponden al modelo base, entrenado sin remuestreo.

\subsection{Evaluación}
Implementamos un evaluador estilo COCO (interpolación de 101 puntos, IoU 0.50:0.05:0.95) con soporte
de cajas \emph{ignoradas}, verificado con pruebas unitarias, porque los análisis por subgrupo
requieren excluir del cómputo las cajas de otros grupos. Reportamos mAP@0.5 y mAP@0.5:0.95 por
especie. Además, a nivel de imagen (\emph{¿hay fauna?}), elegimos el umbral operativo que maximiza
F1 en \texttt{cis\_val} y reportamos precisión, recall, F1, AUC y tasa de falsos positivos en fotos vacías.

\subsection{Robustez ante oclusión y tamaño}
El tamaño relativo es área de la caja / área de la imagen, con cortes en 1\,\% y 10\,\% (pequeño,
mediano, grande). CCT no trae etiquetas de oclusión; definimos un proxy
$s=\tfrac12\,\mathbb{1}[\text{caja truncada por el borde}]+\tfrac12\,\mathrm{rank}(\rho_{\text{bordes}})$,
donde $\rho_{\text{bordes}}$ es la densidad de bordes Canny dentro de la caja (la vegetación
interpuesta fragmenta la silueta), discretizado en terciles. El repositorio genera una muestra
estratificada de 300 cajas de \texttt{trans\_test} para validar el proxy con anotación manual
(kappa ponderado); esa validación no se completó para este informe. Por subgrupo calculamos mAP y
recall al umbral operativo.

\subsection{¿Animal o fondo?}
\textbf{Grad-CAM.} Implementamos Grad-CAM con \emph{hooks} sobre la capa 6 del backbone (C2f, paso 16);
el objetivo es el score de la especie verdadera en la ancla de mayor activación. Medimos la fracción de
energía del mapa dentro de la caja GT y la normalizamos por el área relativa (1 = mapa uniforme).
\textbf{Enmascarado.} Borramos al animal con \emph{inpainting}~\citep{telea2004inpainting} y volvemos a
detectar: si la misma especie sigue apareciendo con confianza sobre el umbral, la decisión depende
del contexto y no del animal.

\subsection{Módulo de monitoreo ecológico}
Sobre todas las fotos de una cámara \emph{trans}, aplicamos reglas: (R1) umbral por especie;
(R2) supresión entre clases con IoU$>0.7$; (R3) foto con fauna si sobrevive una detección;
(R4) agregación por secuencia de disparo, especie por suma de confianzas y conteo MaxN; (R5) fusión de
secuencias de la misma especie separadas por menos de 30 minutos (eventos independientes); (R6)
indicadores: \% de fotos vacías, riqueza, eventos, RAI$=100\cdot$eventos/días-trampa y \% nocturno.
El mismo motor se aplica a la verdad de campo para medir cómo se propagan los errores del detector.

% =====================================================================
\section{Experimentos y resultados}
\subsection{Desempeño por especie y efecto del desbalance}
\tabla{per_species}
\figura{ap_per_species}{0.9}{AP@0.5 por especie en ubicaciones vistas (cis) y nuevas (trans), ordenadas por frecuencia.}
El desbalance es extremo: \emph{opossum} tiene 2\,514 cajas de entrenamiento y \emph{fox} 5 (razón
503:1), y tres cámaras concentran el 53\,\% de las imágenes de entrenamiento. En cámaras vistas, la
frecuencia explica solo parte del desempeño: las clases con más de 1\,000 cajas superan 0.70 de
AP@0.5, pero \emph{skunk} alcanza 0.945 con 214 cajas, mientras que \emph{bird} (560) se queda en 0.614;
la forma y el tamaño del animal pesan tanto como el número de ejemplos. Las cinco especies con menos
cajas promedian 0.499 de AP@0.5 en cis y 0.154 en trans: el desbalance y el cambio de cámara se
suman, aunque el promedio en trans descansa en muy pocas instancias (26 cajas de \emph{rodent}, 15 de
\emph{badger} y una sola de \emph{fox}). \emph{badger} (9 cajas) y \emph{fox} (5) prácticamente no se aprenden, y \emph{deer} no tiene
instancias en \texttt{trans\_test}, por lo que su AP en trans no se define.

La caída cis$\to$trans, sin embargo, no sigue a la frecuencia. \emph{cat} (1\,170 cajas) cae de 0.866
a 0.278 y \emph{squirrel} de 0.707 a 0.275, mientras que \emph{opossum} se mantiene (0.914$\to$0.877) y
\emph{dog} incluso mejora (0.680$\to$0.713); \emph{rodent} colapsa de 0.789 a 0.001, sobre solo 26 cajas. Nuestra hipótesis
es que la generalización depende menos de cuántos ejemplos hay que de en cuántas cámaras distintas
aparece la especie: una especie vista sobre pocos fondos se aprende junto con ellos. \emph{car}, un
objeto grande y rígido, es casi perfecto en ambos casos (0.999 y 0.995).

\tabla{image_level}
Como filtro de fotos vacías, el modelo es muy preciso: al umbral 0.475, el 99\,\% de las fotos que marca
con fauna la tienen, en cámaras vistas y nuevas (Tabla~\ref{tab:imagelevel}). Lo que se degrada es el
recall, de 0.919 a 0.847: en cámaras nuevas, 15 de cada 100 fotos con animal se descartarían como
vacías. La tasa de falsos positivos en fotos vacías sube de 6.4\,\% a 9.6\,\% y el AUC de 0.970 a 0.929,
a pesar de que el modelo nunca vio una foto vacía al entrenar. La precisión, sin embargo, está inflada
porque en estas particiones solo $\sim$7.5\,\% de las fotos están vacías; la cifra que se traslada a un
despliegue real es la tasa de falsos positivos. Con 70\,\% de fotos vacías, como en el conjunto completo,
esa tasa y el recall de trans darían una precisión cercana a 0.79.

\subsection{Generalización a cámaras nuevas}
\figura{map_by_location}{0.38}{mAP@0.5 por cámara: dispersión entre ubicaciones vistas y nuevas.}
El mAP@0.5 baja de 0.720 a 0.455 ($-26.5$ puntos, $-37\,\%$ relativo) y el mAP@0.5:0.95 de 0.484 a
0.310. La dispersión entre cámaras también crece: en las 10 cámaras vistas el mAP@0.5 va de 0.54 a
0.82, y en las 9 nuevas de 0.16 a 0.74 (Figura~\ref{fig:map_by_location}). En la práctica, el
desempeño en un sitio nuevo es difícil de anticipar, que es justamente lo que importa al desplegar.

\tabla{mask_test}
\textbf{Enmascarado.} Es la evidencia más clara (Tabla~\ref{tab:masktest}). Al borrar al animal, el
score de su especie cae en promedio de 0.75 a 0.15 en cámaras vistas y de 0.70 a 0.09 en nuevas: en la
gran mayoría de los casos el modelo sí depende del animal. Pero en el 16.3\,\% de las imágenes de
cámaras vistas la misma especie se sigue detectando sobre el umbral sin el animal, frente al 9.0\,\% en
cámaras nuevas. Esa diferencia es la firma esperada de un atajo de contexto: en las cámaras de
entrenamiento el fondo por sí solo sugiere la especie que suele pasar por ahí, y esa ayuda desaparece
en un sitio nuevo, lo que contribuye a la brecha. Dos precauciones: solo se borra la caja más grande,
de modo que en imágenes con varios individuos de la misma especie la detección puede persistir
legítimamente, y el \emph{inpainting} deja artefactos.

\figura{ejemplos_trans_test}{1.0}{Detecciones en cámaras nuevas (verde = verdad de campo, rojo = predicción,
umbral 0.475). Muestra aleatoria de 12 fotos; los errores típicos son omisiones de animales oscuros u
ocluidos de noche y confusiones de especie con la caja bien localizada.}

\tabla{gradcam}
\textbf{Grad-CAM.} No resultó informativo. La fracción de energía dentro de la caja es similar a la
fracción de la imagen que ocupa la caja (concentración mediana entre 0.67 y 0.97, donde 1 equivale a
un mapa uniforme) en todos los grupos, y no distingue cis de trans (Tabla~\ref{tab:gradcam}); incluso
los errores muestran una concentración algo mayor que los aciertos. Interpretamos que el mapa de la
capa 6 del backbone, calculado sobre una sola ancla, es demasiado difuso para localizar la evidencia
en un detector de una etapa, y no lo usamos como prueba de dependencia del fondo.

\subsection{Robustez ante oclusión y tamaño}
\tabla{robustness}
\textbf{Tamaño.} Es el factor más fuerte. Los animales pequeños ($<$1\,\% de la imagen) tienen recall
0.605 en cis y 0.454 en trans, frente a 0.857 y 0.746 de los medianos, y su mAP@0.5 cae a 0.285 y 0.142.
Los grandes también muestran un mAP bajo (0.489 y 0.202), pero en cis su recall es igual al de los
medianos (0.856): esa caída de AP se debe en buena parte a un artefacto de la métrica por subgrupo,
porque las detecciones sin pareja cuentan como falsos positivos en todos los grupos y pesan más en un
grupo con pocas cajas (1\,621 grandes frente a 10\,494 medianas). Por eso, para comparar grupos de
distinto tamaño usamos el recall. En trans el recall de los grandes sí cae (0.552): animales muy
cercanos a la cámara, a menudo cortados por el borde, son más difíciles sobre un fondo desconocido.

\textbf{Oclusión (proxy).} En cámaras vistas el proxy no separa niveles de desempeño (recall 0.836,
0.853 y 0.817). En cámaras nuevas el recall cae de forma monótona: 0.826 (baja), 0.719 (media) y 0.522
(alta), y el grupo de oclusión alta tiene el menor mAP@0.5 (0.274). La oclusión parece importar
justamente cuando el modelo no puede apoyarse en el conocimiento del fondo. Como el proxy no está
validado con anotación manual, esta conclusión es tentativa.

\subsection{Demostración del módulo de monitoreo}
\tabla{monitor_camera}
\tabla{monitor_species}
\tabla{monitor_periods}
Aplicamos el módulo a la cámara 46 (trans): 6\,276 fotos en 358 días-trampa, entre mayo de 2011 y
junio de 2012. Los totales se acercan a la verdad de campo: 1\,121 eventos independientes frente a
1\,062 ($+5.6\,\%$) y una riqueza de 11 especies frente a 12 (Tabla~\ref{tab:monitorcam}). La
Tabla~\ref{tab:monitorsp} muestra que esa concordancia es una compensación de errores: el modelo inventa
eventos de \emph{car} (30 frente a 0), sobreestima \emph{cat} (82 frente a 5) y \emph{skunk} (82 frente
a 30), subestima \emph{raccoon} (81 frente a 156) y omite por completo \emph{rodent} (14 eventos reales)
y \emph{badger} (1). La lista de especies predicha no es la real, aunque su tamaño casi coincida. En
cambio, las especies frecuentes y bien aprendidas tienen abundancias cercanas a las reales
(\emph{opossum} 310 frente a 327, \emph{bobcat} 208 frente a 225, \emph{rabbit} 123 frente a 142) y
conservan su patrón de actividad (\emph{opossum}: 87\,\% de eventos nocturnos frente a 88\,\%).
El indicador más sesgado es el porcentaje de fotos vacías: 41.4\,\% predicho frente a 11.6\,\% real.
Una parte se explica por las imágenes con especie pero sin caja (15\,\% de \texttt{trans\_test}),
probablemente fotogramas de una secuencia etiquetada en los que el animal ya no es visible, que la
verdad de campo cuenta como \emph{con fauna}. La Tabla~\ref{tab:monitorperiod} muestra el reporte
mensual que el módulo entregaría a un equipo de manejo: la actividad máxima se da en junio y julio de
2011 (177 y 153 eventos), y octubre de 2011 registra solo 2, lo que sugiere un periodo con la cámara
inactiva; por eso el RAI se normaliza por días-trampa. En la práctica, los indicadores de especies
frecuentes son utilizables, pero la lista de especies y la abundancia de las raras requieren revisión
humana.

% =====================================================================
\section{Uso responsable de IA en conservación}
\textbf{Asimetría de costos.} Un falso positivo cuesta segundos de un especialista que descarta
una detección incorrecta; un falso negativo sobre una especie amenazada puede excluirla de los
indicadores con que se decide el manejo de un área (zonificación, patrullaje, presupuesto). Para una
especie con pocas detecciones al año, perder dos o tres eventos cambia su índice de abundancia
relativa de forma material, y el error es invisible porque nadie revisa las fotos ya clasificadas
como vacías. Por eso el umbral no debería ser único: proponemos umbrales por especie más bajos para
especies prioritarias (soportados por la regla R1), una banda de confianza intermedia que se envía a
revisión humana en lugar de descartarse, y auditorías periódicas de una muestra aleatoria de fotos
etiquetadas como vacías para estimar la tasa real de omisión. La métrica que conviene optimizar no es
el mAP global, dominado por especies comunes, sino el recall de las especies de interés con un
presupuesto de revisión acotado. Nuestros resultados lo ilustran: las especies raras pasan de 0.50 a
0.15 de AP@0.5 en cámaras nuevas, y son precisamente las que más suelen importar en conservación. En
la cámara de demostración, \emph{rodent} (14 eventos reales) desaparece por completo del reporte.

\textbf{Despliegue fuera de la región de entrenamiento.} Nuestros propios resultados muestran que el
desempeño cae al pasar a cámaras nuevas dentro de la misma región. Llevar este modelo a la Amazonía o
a los Andes peruanos añade un problema de conjunto cerrado: jaguares, tapires u osos de anteojos no
existen en sus clases y serán asignados con confianza a la especie más parecida o al fondo, con
vegetación, iluminación y tipos de cámara muy distintos. Antes de usar sus salidas para decisiones se
requiere validación local: un lote etiquetado por especialistas en las cámaras del nuevo sitio, recall
por especie con intervalos de confianza, y ajuste fino con datos locales; mientras tanto el modelo
debe operar solo como filtro de vacías con revisión humana de todo lo demás. Por último, las
coordenadas de detecciones de especies amenazadas son información sensible (riesgo de caza furtiva) y
deben protegerse en cualquier reporte o repositorio público.

% =====================================================================
\section{Limitaciones}
(i) El proxy de oclusión es una heurística y no se validó con anotación manual.
(ii) La métrica por subgrupo atribuye las detecciones sin pareja a todos los grupos por igual, lo que
penaliza el AP de los grupos con pocas cajas. (iii) El enmascarado por \emph{inpainting} introduce
artefactos y solo borra la caja más grande de cada imagen. (iv) No entrenamos el experimento con RFS
por el límite de cuota de GPU, y Grad-CAM sobre el backbone resultó demasiado
difuso para ser concluyente. (v) Por presupuesto de cómputo no exploramos YOLOv9 ni resoluciones mayores, que
favorecerían a los animales pequeños. (vi) Las cámaras \emph{trans} pertenecen a la misma región; la
brecha medida es una cota optimista de la brecha entre regiones. (vii) Las imágenes con animal pero
sin caja se excluyen del entrenamiento, lo que puede sesgar la distribución de especies.
(viii) Al usar CCT-20 y no el conjunto completo, entrenamos con menos cámaras y especies, y sin fotos
vacías: incorporar fotos vacías de las mismas ubicaciones desde el conjunto completo es la extensión
más directa, y con más ubicaciones de entrenamiento la brecha \emph{cis/trans} podría reducirse.

\section{Conclusión}
Un YOLOv8s entrenado en 10 cámaras de CCT-20 alcanza 0.720 de mAP@0.5 en esas mismas cámaras pero
0.455 en cámaras nuevas, y las especies raras pasan de 0.50 a 0.15. El test de enmascarado indica que
parte del desempeño en cámaras conocidas se apoya en el contexto de la escena, y los animales pequeños
y, en cámaras nuevas, los ocluidos son los casos más difíciles. En monitoreo, los totales y las
especies frecuentes se estiman bien, pero la lista de especies y la abundancia de las raras no son
confiables sin revisión humana, lo que favorece usar el detector como apoyo y validarlo localmente antes
de desplegarlo en un sitio nuevo.

\paragraph{Reproducibilidad.} Código: \url{https://github.com/mavaldiviezo08/cct-yolo}. Todo el pipeline (descarga,
preparación, entrenamiento, evaluación, figuras y tablas de este informe) se ejecuta con
\texttt{bash run\_all.sh} sobre las imágenes reales de CCT-20. Además, \texttt{bash run\_all.sh -{}-smoke} es una prueba rápida del código (unos minutos, sin GPU ni descarga) con imágenes de prueba generadas; no se usa para ningún resultado de este informe.

{\small\setlength{\bibsep}{2pt}
\bibliographystyle{plainnat}
\bibliography{references}}
\end{document}
